\section{Introduction}\label{sec:intro}

Cyber-physical systems (CPS) couple computation and communication with physical processes through sensing, control and actuation. Attacks on them have physical consequences, and machine-learning anomaly detectors are a widely studied line of defence. A forecaster learns the normal dynamics of sensors, setpoints and actuator commands, and raises an alarm when observations depart from its predictions \citep{kravchik2018cnn,deng2021gdn,fung2024attributions}. An alarm alone, however, does not tell an operator \emph{what} to do. To respond, the operator must know which part of the physical system has been manipulated. The natural unit of a CPS is the feedback control loop: a reference or setpoint (SP), a measured process variable (PV) and a control output (CV) tied together by one controller. It is also the unit in which attacks on industrial CPS are specified \citep{shin2021hai}. Explainable AI (XAI) is the usual answer: post-hoc attribution methods such as SHAP \citep{lundberg2017shap}, LIME \citep{ribeiro2016lime} or integrated gradients \citep{sundararajan2017ig} rank input features by their contribution to the alarm.

Two observations motivate this work. First, the most systematic study to date of attribution for CPS anomaly detection, on industrial control systems (ICS) \citep{fung2024attributions} found that attributions computed at the moment an alarm fires---the only moment that matters operationally---are often wrong. The authors went as far as suggesting that attribution may be better suited to non-real-time workflows. Second, the growing literature that does measure the cost of explanations in cyber-physical settings reports it in absolute terms: seconds per explanation on a GPU \citep{oinonen2026substations}, time per window against a one-second budget \citep{franco2026shats}, or milliseconds on an edge device \citep{ericson2025wfcs}. None of these relates explanation time, or the accuracy of the explanation at the time it is delivered, to what actually bounds the operator's response: the time the physical process can tolerate the attack before reaching an unsafe state. Process-safety engineering formalises this interval as the \emph{process safety time} \citep{iec61511}.

We argue that the right question is not ``how accurate is the explanation?'' nor ``how fast is it?'' but: \emph{does the operator receive a correct localization before the deadline imposed by the process?} This is the question that distinguishes XAI for CPS from XAI for conventional IT intrusion detection, where no physical clock is running. Answering it requires three ingredients the literature lacks:
(i) an evaluation that composes detection delay, explanation latency and the time-varying correctness of the explanation into a single deadline-conditioned quantity;
(ii) public ground truth at the level operators act on, together with physically meaningful deadlines; and
(iii) an attribution method whose accuracy does not collapse during an incident and whose cost is predictable.

\paragraph{Approach} We attribute the detector's anomaly score to the plant's \emph{control loops} with exact Shapley values (\LX). Coalitions are evaluated against an incident-local, last-known-good reference rather than a global background. Grounding the explanation in the control structure of the CPS, rather than in individual input features, makes it both actionable and tractable: with the ten or fewer loops of typical process units, exact computation needs a fixed $2^L$ batched detector evaluations, which yields deterministic, data-independent latency. We combine this with \emph{temporal evidence accumulation} (\TEA): a severity-weighted aggregation of normalised attributions over the successive alarm windows of an incident, with a confidence margin that supports an anytime stopping rule. Within the top-ranked loop, an exact tag-level refinement recovers the manipulated points.

\paragraph{Scope} The method requires only a loop map and black-box access to a detector score, both of which exist for most CPS. Our empirical evidence, however, comes from industrial process CPS, the domain with public attack datasets carrying target ground truth: a real hardware-in-the-loop testbed (HAI) and a process simulator. Other CPS domains (power grids, vehicles, medical devices) are discussed but not evaluated.

\paragraph{Contributions}
\begin{enumerate}
\item \textbf{A deadline-conditioned evaluation framework} for explanations in CPS intrusion detection. We define deadline-conditioned localization DCL@$k(D)$ and an in-time rate against physically derived hazard deadlines. Both account for detection delay, the measured latency of the explainer and the correctness of the explanation actually delivered (\cref{sec:problem}).
\item \textbf{\LX{} and \TEA}: loop-structured exact Shapley attribution with an incident-local reference, and a generic temporal accumulation scheme that applies to any attribution method (\cref{sec:method}).
\item \textbf{Two benchmarks with localization ground truth.} The first is a curated loop- and tag-level ground-truth table for all 50 attacks of the public HAI~21.03 dataset \citep{shin2021hai}, with documented corrections, cross-checked against an independent transcription (49/50 agreement). The second is PCL-Sim, an open closed-loop process simulator with five interacting PI loops. Its \SimNattacks{} attacks follow the HAI taxonomy, and each carries a hazard deadline computed from the true physical state (\cref{sec:setup}).
\item \textbf{A leakage-free, statistically tested comparison} of 12 attribution methods $\times$ 4 detector families $\times$ 3 seeds on both benchmarks. It covers accuracy over time, deadline-conditioned localization, latency and evaluation cost, faithfulness, temporal and perturbation stability, sensor-noise and decoy robustness, and ablations (\cref{sec:results}). A side experiment quantifies how much a random window split inflates detection results on HAI. Such splits appear in several of the XAI-IDS studies we reviewed.
\end{enumerate}

\paragraph{Summary of findings} The results are mixed, and we report them as such.
\begin{itemize}
\item \emph{Timing dominates method choice.} Single-shot attributions of every method degrade after the alarm as the attack propagates through coupled loops. On HAI, \LX{} falls from \lochaiLXalarmZerosinglehrOne{} top-1 accuracy at the alarm to \lochaiLXalarmOneEightZerosinglehrOne{} after 180\,s. \TEA{} recovers most of this loss for every method (\lochaiLXalarmOneEightZeroteahrOne{} for \LX) and makes the top-1 output temporally stable.
\item \emph{Control-loop structure is essential.} Exact Shapley values over correlation clusters or random partitions of the same size are far less accurate than over control loops.
\item \emph{\LX{} is competitive, not dominant.} It is the most accurate black-box method at the alarm on HAI. With \TEA, however, it is statistically tied with integrated gradients and other strong baselines, and on PCL-Sim integrated gradients and a training-background variant are slightly better. Its distinguishing properties are a fixed, data-independent cost and gradient-free determinism.
\item \emph{Most deadline misses are detection misses.} On hazardous simulated attacks, \intimeFailUndeliveredPct\,\% of the deadline failures of \LX+\TEA{} occur because the attack is detected too late or not at all.
\item \emph{Leakage matters.} On HAI, a random window split inflates supervised detection F1 from \leakBlockedFOne{} to \leakRandomFOne.
\end{itemize}

